\documentclass[12pt,a4paper,oneside]{article}
\usepackage[brazil]{babel}
\usepackage[utf8]{inputenc}
\usepackage{olymp}

\setcounter{problem}{3}

\begin{document}

\begin{problem}{Consumo de combustível}{consumo}{2}

\small
Uma empresa de logística possui uma frota de caminhões que utiliza para transportar cargas por todo o Brasil. Um consultor contratado pela empresa mostrou ao gerente que alguns caminhões poderiam estar gastando mais combustível que o normal, elevando o custo da entrega e diminuindo o lucro.

O gerente precisa aferir a média de consumo de combustível de cada veículo para identificar os veículos problemáticos. Foi então definido que todo motorista deve receber o caminhão com o tanque cheio e, durante a viagem, deve anotar a quilometragem percorrida em cada trecho e os possíveis abastecimentos realizados. Ao final de uma viagem, ao entregar o caminhão, o tanque deve necessariamente ser completado, não esquecendo de anotar a quantidade de combustível.

Dessa forma, o calculo do consumo médio é simplificado, sendo este a distância total percorrida dividida pela quantidade total de combustível abastecido.

Escreva um programa para ajudar a empresa a calcular o consumo médio de combustível de um caminhão em uma determinada viagem.

%\Notes
%
%Observações, se tiver.

\InputFile

A entrada é composta por várias linhas. Cada linha inicia com uma letra $L$ indicando o tipo da entrada e um valor real $V$.

\begin{itemize}
\vspace{-9pt}
    \item Para $L$ = `\texttt{R}', $V$ representa o número de quilômetros rodados em um trecho;
\vspace{-9pt}
    \item Para $L$ = `\texttt{A}', $V$ representa a quantidade de combustível em litros abastecida.
\end{itemize}

\vspace{-9pt}
A entrada se encerra com a letra `\texttt{X}', quando então deve-se calcular e imprimir a média de consumo. Restrições: $0 < V \leq 5000$ e a entrada não começa com `\texttt{X}'.

\OutputFile

Seu programa deve imprimir um única linha, contendo a média de consumo de combustível da viagem informada na entrada (nesta questão, em Km/l). A saída deve ser formatada para duas casas decimais.

% \Notes
% xxx

\Example

\begin{example}
\exmp{
R 100
A 20
R 300
R 200
A 160
X
}{
3.33
}%
\end{example}

\begin{example}
\exmp{
R 200
A 26.98
R 36.15
R 69.54
R 152
A 75.2
X
}{
4.48
}%
\end{example}

\begin{example}
\exmp{
R 200
A 20
X
}{
10.00
}%
\end{example}


%Comentários sobre o exemplo, se necessário.

\end{problem}

\end{document}
